\documentclass[11pt,a4paper]{article}

\usepackage[T1]{fontenc}
\usepackage[utf8]{inputenc}
\usepackage{lmodern}
\usepackage[margin=2.5cm]{geometry}
\usepackage{amsmath,amssymb}
\usepackage{booktabs}
\usepackage{graphicx}
\usepackage{xcolor}
\usepackage[round,authoryear]{natbib}
\usepackage[colorlinks=true,linkcolor=blue,citecolor=blue,urlcolor=blue]{hyperref}

% Notation shared by all sections (follows the portfolio construction course)
\newcommand{\vect}[1]{\boldsymbol{#1}}
\newcommand{\w}{\vect{w}}
\newcommand{\muv}{\vect{\mu}}
\newcommand{\one}{\vect{1}}
\newcommand{\Sig}{\boldsymbol{\Sigma}}
\newcommand{\E}{\mathbb{E}}
\newcommand{\V}{\mathbb{V}}
\newcommand{\Cov}{\operatorname{Cov}}
\newcommand{\todo}[1]{\textcolor{red}{[TODO: #1]}}

% Decision log entry: \decision{ID}{date}{title}{decision}{alternatives}{rationale}
\newcommand{\decision}[6]{%
  \subsection*{#1 -- #3}
  \addcontentsline{toc}{subsection}{#1 -- #3}
  \begin{description}
    \item[Date] #2
    \item[Decision] #4
    \item[Alternatives] #5
    \item[Rationale] #6
  \end{description}}


\title{FX Portfolio Lab:\\From Portfolio Construction to FX Structuring}
\author{Tom Rousselet\\EDHEC Business School\\\href{mailto:tom.rousselet@edhec.com}{\texttt{tom.rousselet@edhec.com}}}
\date{Working paper -- draft of \today}

\begin{document}
\maketitle

\begin{abstract}
\noindent
This working paper documents a research project that studies the foreign exchange (FX) market from three angles: the portfolio manager who allocates across currencies, the investor who hedges the currency risk of foreign assets, and the structurer who designs hedging products for corporate clients.
We first build a valuation and backtesting engine whose correctness is verified against analytical results.
We then construct G10 currency portfolios from carry, momentum and value signals and combine them with minimum-variance, mean-variance and risk-parity allocations, paying particular attention to estimation error and transaction costs.
Next, we derive the optimal hedge ratio of a euro-based investor.
Finally, we price and compare the structured products used to hedge FX exposures.
Every design decision is recorded together with the alternatives considered and the reason for the choice (Appendix~\ref{app:decisions}).
\end{abstract}

\tableofcontents
\bigskip

\section{Introduction}\label{sec:intro}

The FX market is the largest financial market in the world, with an average daily turnover of about 7.5 trillion US dollars in April 2022 \citep{bis2022}.
Banks intervene in it in three main roles: they take and manage risk on behalf of investors, they help clients hedge their exposures, and they design structured solutions for corporate treasurers.
This project reproduces, at a small scale and with public data, the analytical work behind each of these roles.

The project is organised around five research questions.

\begin{enumerate}
  \item \textbf{Engine.} How can a currency portfolio be valued correctly, including forward positions, rebalancing and transaction costs? (Section~\ref{sec:engine})
  \item \textbf{Currency portfolios.} Do the carry, momentum and value strategies documented in the literature \citep{lustig2011,menkhoff2012carry,menkhoff2012momentum,asness2013} survive realistic transaction costs, and which allocation method combines them best out of sample? (Section~\ref{sec:portfolios})
  \item \textbf{The euro.} How much of the monthly moves of the euro is explained by euro-area rates and by sovereign risk? (Section~\ref{sec:euro})
  \item \textbf{Hedging.} Which fraction of the currency exposure of foreign equities and bonds should a euro-based investor hedge to minimise risk? \citep{campbell2010} (Section~\ref{sec:hedging})
  \item \textbf{Structuring.} How do forwards, options and structured products (collars, participating forwards, target redemption forwards) compare for a corporate hedger, in terms of cost and distribution of outcomes? (Section~\ref{sec:structuring})
\end{enumerate}

The analytical framework of the first two questions is that of portfolio construction theory: valuation of rebalanced portfolios, first two moments of a portfolio return, minimum-variance and mean-variance optimisation \citep{markowitz1952,milhau2026}.
A recurring theme is that the inputs of these optimisations are estimated, and that naive mean-variance portfolios are dominated out of sample by simple rules \citep{demiguel2009}.

\paragraph{Methodological principles.}
Three principles are applied throughout.
(i) \emph{No look-ahead bias}: every quantity used to form a portfolio at date $t$ is computed from data available at $t$.
(ii) \emph{Verification}: every piece of code is checked against an analytical result or a published number before it is used (Appendix~\ref{app:software}).
(iii) \emph{Traceability}: every choice that could change the results is recorded in the decision log (Appendix~\ref{app:decisions}).

\section{Framework and notation}\label{sec:framework}

\subsection{Portfolio valuation}

A portfolio invested in $N$ constituents holds an amount $\theta_{t,i}$ in constituent $i$ at date $t$.
Its value and weights are
\begin{equation}
  V_t = \sum_{i=1}^N \theta_{t,i}, \qquad w_{t,i} = \frac{\theta_{t,i}}{V_t}.
\end{equation}
Between two rebalancing dates the portfolio is buy-and-hold: the number of units of each constituent is constant, so each amount grows with the constituent's own return,
\begin{equation}\label{eq:amount-update}
  \theta_{[t+1]^-,i} = \theta_{t,i}\,(1 + r_{t,t+1,i}),
\end{equation}
and the portfolio return is the average of the constituents' returns weighted by the \emph{initial} weights,
\begin{equation}\label{eq:fundamental}
  r_{t,t+1,p} = \sum_{i=1}^N w_{t,i}\, r_{t,t+1,i} = \w_t' \vect{r}_{t,t+1}.
\end{equation}
On a rebalancing date the amounts are reset by the allocation rule; for a fixed-mix rule with target weights $w_i$, $\theta_{t+1,i} = w_i V_{t+1}$.
The amount traded, $\sum_i |\theta_{t+1,i} - \theta_{[t+1]^-,i}|$, determines the transaction costs.

\subsection{Moments and optimal portfolios}

With $\muv$ the vector of expected returns and $\Sig$ the covariance matrix of the constituents,
\begin{equation}
  \mu_p = \w'\muv, \qquad \sigma_p^2 = \w'\Sig\w.
\end{equation}
The global minimum variance (GMV) portfolio and the maximum return-to-risk (MRR) portfolio are
\begin{equation}
  \w_{\mathrm{gmv}} = \frac{\Sig^{-1}\one}{\one'\Sig^{-1}\one}, \qquad
  \w_{\mathrm{mrr}} = \frac{\Sig^{-1}\muv}{\one'\Sig^{-1}\muv},
\end{equation}
and, by the two-fund separation theorem, every efficient portfolio is a combination of the two \citep{milhau2026}.

\subsection{Currency excess returns}

Let $S_t$ be the spot exchange rate, in US dollars per unit of foreign currency, and $i_t$, $i_t^*$ the domestic (USD) and foreign interest rates over $(t,t+1)$.
A US investor who borrows one dollar at $i_t$ and invests it in a foreign deposit earns the excess return
\begin{equation}\label{eq:rx-deposit}
  rx_{t+1} = (1 + i_t^*)\frac{S_{t+1}}{S_t} - (1 + i_t)
  \approx \underbrace{(i_t^* - i_t)}_{\text{carry}} + \underbrace{\frac{S_{t+1} - S_t}{S_t}}_{\text{spot change}}.
\end{equation}
The same exposure is obtained with a forward contract.
Under covered interest parity (CIP) the forward rate is $F_t = S_t (1 + i_t)/(1 + i_t^*)$, and a long forward position on $1/F_t$ units of foreign currency yields
\begin{equation}\label{eq:rx-forward}
  \frac{S_{t+1} - F_t}{F_t} = \frac{rx_{t+1}}{1 + i_t}.
\end{equation}
The approximation in \eqref{eq:rx-deposit} shows the economic content of a currency position: the investor earns the interest rate differential and bears the risk of a depreciation of the foreign currency.
Uncovered interest parity predicts that the two terms offset on average; the empirical failure of this prediction is the source of the carry trade premium \citep{lustig2011}.
Since 2008, CIP no longer holds exactly \citep{du2018}; the consequences for our data are discussed in Section~\ref{sec:data}.

\section{Data}\label{sec:data}

All series are downloaded from FRED, the database of the Federal Reserve Bank of St.~Louis, by the script \texttt{scripts/download\_data.py} (Decision~D7).
Table~\ref{tab:data} lists them.

\begin{table}[h]
\centering
\small
\begin{tabular}{@{}lllll@{}}
\toprule
Currency & Spot series & Spot available & Rate series & Rate available \\
\midrule
USD & -- & -- & \texttt{IR3TIB01USM156N} & 1964-06 -- 2026-09 \\
EUR & \texttt{DEXUSEU} & 1999-01 -- 2026-10 & \texttt{IR3TIB01EZM156N} & 1994-01 -- 2026-01 \\
GBP & \texttt{DEXUSUK} & 1971-01 -- 2026-10 & \texttt{IR3TIB01GBM156N} & 1957-01 -- 2026-01 \\
AUD & \texttt{DEXUSAL} & 1971-01 -- 2026-10 & \texttt{IR3TIB01AUM156N} & 1968-01 -- 2026-08 \\
NZD & \texttt{DEXUSNZ} & 1971-01 -- 2026-10 & \texttt{IR3TIB01NZM156N} & 1973-12 -- 2026-09 \\
JPY & \texttt{DEXJPUS} & 1971-01 -- 2026-10 & \texttt{IR3TIB01JPM156N} & 2002-04 -- 2026-08 \\
CHF & \texttt{DEXSZUS} & 1971-01 -- 2026-10 & \texttt{IR3TIB01CHM156N} & 1999-07 -- 2026-09 \\
CAD & \texttt{DEXCAUS} & 1971-01 -- 2026-10 & \texttt{IR3TIB01CAM156N} & 1956-01 -- 2026-09 \\
NOK & \texttt{DEXNOUS} & 1971-01 -- 2026-10 & \texttt{IR3TIB01NOM156N} & 1979-01 -- 2026-09 \\
SEK & \texttt{DEXSDUS} & 1971-01 -- 2026-10 & \texttt{IR3TIB01SEM156N} & 1982-01 -- 2026-09 \\
\bottomrule
\end{tabular}
\caption{Data series (downloaded on 19 October 2026). Spot: Federal Reserve H.10, daily noon buying rates in New York. Rates: OECD Main Economic Indicators, 3-month interbank rates, monthly averages, \% per year.}
\label{tab:data}
\end{table}

\paragraph{Quotation convention.}
FRED quotes the euro, the pound and the Australian and New Zealand dollars in US dollars per unit of foreign currency, and the other currencies in foreign currency per US dollar.
All spot rates are converted to US dollars per unit of foreign currency, so that an increase of $S_t$ is always an appreciation of the foreign currency, consistently with equation~\eqref{eq:rx-deposit}.

\paragraph{Sample.}
The common sample of spot and interest rates starts in April 2002 (Japanese rates) and ends in January 2026 (euro area and UK rates).
\todo{decide between a balanced panel 2002--2026 and an unbalanced panel that uses each currency as soon as it is available.}

\paragraph{Limitations.}
\begin{itemize}
  \item \textbf{Forward rates} are not freely available; they are implied by CIP (equation~\ref{eq:rx-forward}). Since 2008, CIP deviations of the order of a few tens of basis points per year exist \citep{du2018}.
  \item \textbf{Interest rates} are monthly averages of interbank rates, not the rate observed on the day a position is opened.
  \item \textbf{Transaction costs}: bid--ask spreads are not in the data and must be modelled \todo{reference and figures}.
\end{itemize}

\section{Valuation and backtesting engine}\label{sec:engine}

The engine values a portfolio by propagating the dollar amounts held in each constituent, following the general method of \citet{milhau2026}: amounts are updated with equation~\eqref{eq:amount-update} between rebalancing dates and reset by the allocation rule on rebalancing dates.
Working with amounts rather than weights has two advantages.
First, it handles any rebalancing frequency with the same code (daily rebalancing and buy-and-hold are the two limiting cases).
Second, it gives the amount traded at each date directly, from which transaction costs are computed (Decision~D4).

\paragraph{Verification.}
The engine reproduces the worked examples of the course: the valuation of an equally weighted portfolio of five US stocks between 2 and 5 January 2015 (final value \$978.21, \$6.37 traded to return to equal weights), the drift of weights in a buy-and-hold portfolio, and the compounding of returns.
These examples are encoded as automated tests (Appendix~\ref{app:software}).

\todo{Phase 1: forward positions, transaction costs, performance and risk measures.}

\section{Currency portfolios}\label{sec:portfolios}

\todo{Phase 2.}

\begin{itemize}
  \item Signals: carry (interest rate differential), momentum (past excess returns), value (deviation from purchasing power parity).
  \item Portfolio formation: sorting into long-short portfolios versus optimisation.
  \item Allocation across strategies: equal weights, GMV, mean-variance, risk parity; covariance estimation with shrinkage \citep{ledoit2004}.
  \item Crash risk of the carry trade \citep{brunnermeier2008}.
\end{itemize}

\section{Optimal currency hedging for a euro-based investor}\label{sec:hedging}

\todo{Phase 3.}

The minimum-variance hedge ratio is a direct application of the two-asset minimum-variance weight $w_2^* = (\sigma_1^2 - \sigma_{12})/\sigma_{2/1}^2$ \citep{milhau2026}, where the two assets are the unhedged foreign asset and the currency forward.
The empirical question is whether some currencies, such as the US dollar, act as natural hedges for equity risk \citep{campbell2010}.

\section{FX structuring}\label{sec:structuring}

\todo{Phase 4.}

\begin{itemize}
  \item Option pricing: \citet{garman1983}; volatility smile from at-the-money, risk-reversal and butterfly quotes.
  \item Products: forward, vanilla option, collar, participating forward, target redemption forward (Monte Carlo).
  \item Comparison for a corporate hedger: upfront cost, distribution of the hedged rate, worst-case outcomes.
\end{itemize}

\section{Conclusion and limitations}\label{sec:conclusion}

\todo{Written at the end of the project.}

Known limitations at this stage:
\begin{itemize}
  \item Forward rates are implied by covered interest parity rather than observed.
  \item Transaction costs are modelled, not observed.
\end{itemize}


\bibliographystyle{plainnat}
\bibliography{references}

\appendix
\section{Decision log}\label{app:decisions}

Each entry records a choice that could affect the results or the usability of the project, the alternatives that were considered and the reason for the choice.
Entries are never deleted: a decision that is reversed is superseded by a new entry that refers to it.

\decision{D1}{18 October 2026}{Language of the project}
{The paper, the code and its comments are written in English.}
{French.}
{English is the working language of trading floors and of open-source code, and the project is meant to be shown to recruiters, including outside France.}

\decision{D2}{18 October 2026}{Code organised as an installable Python package}
{The code lives in a package (\texttt{src/fxlab}) installed in a virtual environment; notebooks only import and use it.}
{(a) Code written directly in Jupyter notebooks. (b) Loose scripts.}
{Code in notebooks cannot be tested or reused, and is easily run out of order. A package separates the library (tested, reused by every phase) from the analysis (notebooks). The \texttt{src} layout guarantees that the tests run against the installed package, not against stray files.}

\decision{D3}{18 October 2026}{Course examples used as automated tests}
{Every worked example of the course that the code can reproduce is encoded as a test (\texttt{pytest}), with a tolerance that reflects the rounding of the published inputs.}
{Checking results by eye in notebooks.}
{The course examples have known answers, so they are an independent specification of what the code must do. Automated tests are rerun after every change, which prevents an improvement in one place from silently breaking another.}

\decision{D4}{18 October 2026}{Valuation by propagation of amounts}
{Portfolios are valued by propagating the amount held in each constituent and resetting it on rebalancing dates (course, section~1.8).}
{Compounding the portfolio return $\w'\vect{r}$ with fixed weights.}
{Compounding with fixed weights is only valid when the portfolio is rebalanced at every date. Propagating amounts handles any rebalancing frequency and gives the amount traded, which is needed for transaction costs.}

\decision{D5}{18 October 2026}{Free, official data sources}
{Exchange rates and interest rates are taken from official publishers (Federal Reserve, OECD, ECB) through FRED and the ECB data portal; forward rates are implied by CIP.}
{(a) Bloomberg or Refinitiv: no access. (b) Yahoo Finance: unofficial interface, no documentation of the series, known data errors.}
{Official sources are documented, stable and allow anyone to reproduce the results. The cost is that forward rates are not observed; this is acknowledged as a limitation (Section~\ref{sec:data}).}

\decision{D6}{18 October 2026}{Working paper in \LaTeX{}, versioned with the code}
{The documentation is a \LaTeX{} article stored in the same repository as the code.}
{(a) A word processor document. (b) Notebooks only.}
{\LaTeX{} is the standard for scientific documents with equations. Keeping it in the repository ties each version of the text to the version of the code that produced its results.}

\decision{D7}{19 October 2026}{Data downloaded by a script from the FRED CSV interface}
{A script downloads each series from its public FRED address and caches the raw file in \texttt{data/raw/}; the files are read and cleaned by the \texttt{fxlab.data} module. Spot rates are converted to US dollars per unit of foreign currency.}
{(a) Manual download from the website. (b) The FRED API, which requires a personal key. (c) A third-party package such as \texttt{pandas-datareader}.}
{A script makes the dataset reproducible with one command and can refresh it at any time. The public CSV interface needs no key and only the Python standard library, so there is no dependency that could break. A single quotation convention avoids sign errors on half of the currencies.}

\decision{D8}{19 October 2026}{Monthly frequency, month-end sampling}
{Currency strategies are formed and evaluated monthly, using the last available spot rate of each month.}
{(a) Daily frequency. (b) Weekly frequency. (c) Monthly averages of spot rates.}
{Interest rates are only available monthly, so a higher frequency would reuse the same rate for a whole month. Monthly rebalancing is the standard of the currency literature \citep{lustig2011,menkhoff2012carry}, which makes the results comparable, and it keeps turnover and transaction costs low. Month-end rates, not averages, are used because a position is opened and closed at a given rate, not at an average.}

\decision{D9}{19 October 2026}{Unbalanced panel from January 1999}
{The sample starts in January 1999; each currency enters as soon as its spot rate and interest rate are both available. A balanced panel (2002--2026) is used as a robustness check.}
{(a) Balanced panel from May 2002 only. (b) Start in the 1980s without the euro.}
{Starting with the euro gives a homogeneous G10 universe. An unbalanced panel uses three more years of data without inventing any, as in \citet{lustig2011}; the robustness check verifies that the conclusions do not depend on this choice.}

\decision{D10}{19 October 2026}{US dollar as base currency for the currency portfolios}
{The excess returns of Phase 2 are measured from the point of view of a US investor; the euro-based investor is the subject of Phase 3.}
{Euro as base currency from the start.}
{The data are quoted against the US dollar and almost all of the literature uses a US investor, so the results can be compared with published ones. Excess returns of long-short portfolios do not depend much on the base currency, since the base currency cancels out of a long-short position.}

\decision{D11}{19 October 2026}{Excess return defined as the simple return of a one-month forward}
{$rx_{t+1} = S_{t+1}/F_t - 1$ with $F_t$ implied by CIP and the monthly rate equal to the annual rate divided by 12 (equations~\ref{eq:fwd-monthly}--\ref{eq:rx-monthly}).}
{(a) Log excess return $\ln S_{t+1} - \ln F_t$, common in the literature. (b) Deposit version \eqref{eq:rx-deposit}.}
{The portfolio formula $r_p = \w'\vect{r}$ of the course holds for simple returns, not for log returns, so simple returns can be aggregated across currencies without approximation. The forward is the instrument actually traded by an FX desk for this exposure, and it requires no capital beyond margin. The deposit version differs only by the factor $1+i_t$ (equation~\ref{eq:rx-forward}).}

\decision{D12}{19 October 2026}{Continuous integration and MIT licence}
{The test suite runs automatically on GitHub Actions at every push, on Python 3.11 and 3.13; the code is published under the MIT licence.}
{(a) Running the tests locally only. (b) No licence, or a restrictive one.}
{Tests run on a clean machine catch errors that a local environment hides (a missing dependency, a file that only exists locally), and the badge shows any reader that the suite passes. Without a licence, published code cannot legally be reused; MIT is the most common licence for demonstration projects and imposes only attribution.}

\decision{D13}{19 October 2026}{Compiled paper versioned in the repository}
{The compiled PDF is copied to \texttt{paper/fx-portfolio-lab.pdf} by \texttt{make} and committed with each update of the paper; LaTeX build files remain ignored.}
{(a) PDF attached to GitHub releases only. (b) PDF compiled by the continuous integration. (c) Sources only.}
{A reader who is not going to compile \LaTeX{} must be able to open the paper in one click from the README. Committing the PDF is the simplest way to guarantee it and keeps each version of the PDF next to the code that produced it. Compiling in the CI would require installing a TeX distribution on every run, which is slow for a document that changes only a few times per phase.}

\decision{D14}{19 October 2026}{Transaction costs: roll plus change of position}
{The notional kept from one month to the next pays the swap half-spread (roll); the notional traded pays the outright half-spread (equation~\ref{eq:cost}).}
{A single half-spread on the notional traded only.}
{Holding a forward position requires rolling it every month with an FX swap, which is how an FX desk maintains the exposure. Ignoring the roll would understate the cost of strategies that hold their positions; on the dollar factor, the roll accounts for most of the cost (Section~\ref{sec:dollar}).}

\decision{D15}{19 October 2026}{Cost level: stated assumption, sensitivity and break-even}
{The base-case half-spreads are an explicit assumption by liquidity tier (Section~\ref{sec:costs}); results are reported for the base case scaled by 0, 1, 2 and 5, with the break-even cost scale of each strategy.}
{(a) Figures presented as estimates without a public data source. (b) A single cost for all currencies.}
{Bid--ask quotes are not freely available, so any single figure would be an assumption; presenting it as an estimate would overstate its precision. The break-even scale answers the question that matters (up to which cost does a strategy remain profitable?) without depending on the assumption, and the tiers keep the ranking of liquidity between currencies.}

\decision{D16}{20 October 2026}{Interbank rates in covered interest parity, with a risk-free-rate robustness check}
{Forwards are implied from 3-month interbank rates over the whole sample; from November 2019, euro and sterling excess returns are recomputed with the overnight risk-free rates (SOFR, ESTR, SONIA) as a robustness check (Section~\ref{sec:rates-choice}).}
{(a) Risk-free rates only: they do not exist before 2018--2019 and not for every currency in free data. (b) Interbank rates without discussion.}
{Interbank rates contain a bank credit premium and a term component that the rates of a collateralised FX swap do not contain, a point raised by practitioners. But they were the market reference before the benchmark reform, they cover the whole sample and all currencies, and the robustness check shows that switching to risk-free rates changes monthly excess returns by about 1 basis point. The choice is therefore kept and its limits are stated.}

\decision{D17}{20 October 2026}{Extension of the euro and sterling rates after January 2026}
{The euro rate is continued with the 3-month Euribor published by the ECB; the sterling rate with the monthly average of SONIA, shifted by the mean gap over the last twelve months of overlap.}
{(a) Let the two currencies leave the panel in February 2026. (b) Splice SONIA without a shift.}
{The OECD euro series is the Euribor, so the ECB series continues it exactly. For sterling no 3-month interbank rate exists after LIBOR; a constant shift removes the jump at the junction ($-0.08$ point) without changing the monthly variation of the rate. Keeping the euro and sterling in the panel avoids changing the investment universe at the end of the sample.}

\decision{D18}{20 October 2026}{Euro drivers measured with monthly averages and Newey--West standard errors}
{The regressions of Section~\ref{sec:euro} use changes of monthly averages for the exchange rates as well as for rates and yields, and Newey--West standard errors with the usual lag rule.}
{(a) Month-end exchange rates against monthly-average yields. (b) Classical OLS standard errors.}
{The OECD rates and yields are only available as monthly averages; mixing them with month-end exchange rates would misalign the timing of the two sides. Averaging induces autocorrelation in monthly changes, and financial data are heteroskedastic, so classical standard errors would overstate precision; Newey--West standard errors are robust to both.}

\decision{D19}{20 October 2026}{Scope: what is not added from the practitioner advice}
{ECB meeting pricing, sovereign debt auctions and the carry and roll-down of bonds are not part of the project; they are followed as market monitoring outside it.}
{(a) An event study of ECB meeting days. (b) A rates module for bond carry, roll-down and repo.}
{Measuring a monetary policy surprise requires the market pricing of each meeting (dated overnight-index swaps), which is not freely available historically, and a complete list of meeting dates that would have to be compiled by hand; auctions have no free historical database. Bond carry and roll-down belong to a rates desk: adding them would blur a project whose subject is currencies. The euro is studied instead through the two channels of Section~\ref{sec:euro}, which public data can measure.}

\subsection*{Open decisions}
\addcontentsline{toc}{subsection}{Open decisions}
None at this stage. The decisions of Phase 2 (signal construction, portfolio formation, estimation of the covariance matrix) will be added before the corresponding code is written.

\section{Software architecture and verification}\label{app:software}

\begin{table}[h]
\centering
\small
\begin{tabular}{@{}ll@{}}
\toprule
Module & Content \\
\midrule
\texttt{fxlab.portfolio} & Returns, compounding, valuation of buy-and-hold and rebalanced portfolios \\
\texttt{fxlab.data} & Download and loading of FRED and ECB series, quotation convention, splicing \\
\texttt{fxlab.fx\_returns} & Month-end sampling, CIP forwards, excess returns, carry \\
\texttt{fxlab.costs} & Cost model: roll and outright half-spreads, base case and scaling \\
\texttt{fxlab.forwards} & Backtest of portfolios of rolled forwards, equal weights \\
\texttt{fxlab.metrics} & Performance and risk measures, break-even cost scale \\
\texttt{fxlab.regression} & OLS with Newey--West standard errors, rolling coefficients \\
\texttt{fxlab.euro} & Monthly changes of the euro and of its rates and risk drivers \\
\bottomrule
\end{tabular}
\caption{Modules of the \texttt{fxlab} package.}
\end{table}

\begin{table}[h]
\centering
\footnotesize
\begin{tabular}{@{}llp{5.6cm}@{}}
\toprule
Test & Source & Result checked \\
\midrule
\texttt{test\_table1\_buy\_and\_hold} & Course, Table 1 & $V_1 = \$978.21$, return $-2.18\%$ \\
\texttt{test\_table1\_rebalancing\_\ldots} & Course, section 1.4 & \$6.37 traded \\
\texttt{test\_midterm\_q1\_amount\_update} & Midterm 2024--25, Q1 & \$36.29 \\
\texttt{test\_midterm\_q4\_weight\_drift} & Midterm 2024--25, Q4 & weight 74.9\% \\
\texttt{test\_midterm\_q7\_compounding} & Midterm 2024--25, Q7 & $+0.8\%$ \\
\texttt{test\_single\_rebalance\_\ldots} & Course, section 1.8 & equals buy-and-hold \\
\texttt{test\_spot\_converted\_\ldots} & FRED conventions & JPY per USD inverted \\
\texttt{test\_missing\_values\_\ldots} & FRED format & missing values are NaN \\
\texttt{test\_rates\_in\_decimal} & FRED format & 5\% read as 0.05 \\
\texttt{test\_forward\_by\_cip} & Eq.~\eqref{eq:fwd-monthly} & forward computed by hand \\
\texttt{test\_excess\_return\_by\_hand} & Eq.~\eqref{eq:rx-monthly} & excess return computed by hand \\
\texttt{test\_equal\_rates\_\ldots} & Eq.~\eqref{eq:rx-monthly} & no carry: $rx$ = spot return \\
\texttt{test\_carry\_\ldots} & Section~\ref{sec:rx} & carry $\approx (i^*-i)/12$ \\
\texttt{test\_no\_look\_ahead} & Timing rule & $rx_{t+1}$ unaffected by later data \\
\texttt{test\_month\_end\_spot\_\ldots} & Sampling rule & last available day used \\
\texttt{test\_incomplete\_last\_month\_\ldots} & Sampling rule & partial month dropped \\
\texttt{test\_rates\_dated\_\ldots} & Timing rule & rate dated at month end \\
\midrule
\texttt{test\_costs.py} (10 tests) & Eq.~\eqref{eq:cost} & opening, roll, resize, sign change, scaling, by hand \\
\texttt{test\_single\_long\_\ldots} & Eq.~\eqref{eq:fwd-value} & one forward earns $rx$ \\
\texttt{test\_long\_short\_\ldots} & Eq.~\eqref{eq:fwd-value} & weights summing to 0: $r_p = \w'\vect{rx}$ \\
\texttt{test\_collateral\_\ldots} & Eq.~\eqref{eq:fwd-value} & collateral earns $i_t/12$ \\
\texttt{test\_costs\_by\_hand} & Eqs.~\eqref{eq:fwd-value}--\eqref{eq:cost} & costs of two months by hand \\
\texttt{test\_turnover\_of\_a\_flip} & Section~\ref{sec:metrics} & long to short trades twice the notional \\
\texttt{test\_no\_look\_ahead} (forwards) & Timing rule & weights at $t$ only affect month $t+1$ \\
\texttt{test\_position\_without\_\ldots} & Section~\ref{sec:fwd-backtest} & position without a return rejected \\
\texttt{test\_metrics.py} (7 tests) & Section~\ref{sec:metrics} & mean, vol, Sharpe, $t$-stat, drawdown, skewness by hand \\
\midrule
\texttt{test\_data.py}, new tests (5) & Section~\ref{sec:rates-choice} & ECB file format, monthly averages, splice without jump, overlap required \\
\texttt{test\_regression.py} (8 tests) & Section~\ref{sec:euro} & coefficients recovered, White formula, HAC wider under autocorrelation, lag rule, rolling, drivers by hand \\
\bottomrule
\end{tabular}
\caption{Verification tests.}
\end{table}


\end{document}
